% Preámbulo común del solucionario (lo usan el documento maestro y las pruebas por tema).
% \solroot debe definirse ANTES de % Preámbulo común del solucionario (lo usan el documento maestro y las pruebas por tema).
% \solroot debe definirse ANTES de % Preámbulo común del solucionario (lo usan el documento maestro y las pruebas por tema).
% \solroot debe definirse ANTES de % Preámbulo común del solucionario (lo usan el documento maestro y las pruebas por tema).
% \solroot debe definirse ANTES de \input{preambulo}: "solucionario/" en el maestro, vacío en las pruebas.
\usepackage[spanish,es-noshorthands,es-nodecimaldot,es-tabla]{babel}
\usepackage{fontspec}
\usepackage{amsmath,amssymb,mathtools}
\usepackage{booktabs,array}
\usepackage{xcolor}
\usepackage{enumitem}
\usepackage{graphicx}
\usepackage{float}
\usepackage{siunitx}
\usepackage{fancyvrb}
\usepackage[letterpaper,margin=2.2cm]{geometry}
\usepackage{titlesec}
\usepackage{fancyhdr}
\usepackage{hyperref}
\hypersetup{colorlinks=true,linkcolor=black,urlcolor=black}
\sisetup{output-decimal-marker={.}, per-mode=symbol, range-phrase={--}}
\graphicspath{{\solroot figuras/}}
\fvset{fontsize=\scriptsize, frame=single, numbers=left, numbersep=5pt}

\definecolor{tema}{RGB}{20,60,120}
\definecolor{gris}{RGB}{90,90,90}
\definecolor{caja}{RGB}{235,241,250}
\titleformat{\section}{\Large\bfseries\color{tema}}{\thesection}{0.6em}{}
\titleformat{\subsection}{\large\bfseries\color{tema}}{}{0em}{}
\setlist{itemsep=1pt,topsep=2pt}
\setlength{\parindent}{0pt}
\setlength{\parskip}{0.6ex}
\pagestyle{fancy}\fancyhf{}
\lhead{\small Óptica No Lineal 2026-II — Solucionario del taller de repaso, Examen 1}
\rhead{\small p.~\thepage}

% Encabezado de cada solución: \sol{número}{título}{fuente}
\newcommand{\sol}[3]{\subsection{Problema #1. #2}\noindent{\small\color{gris}#3}\par}
% Inciso: \inc{a}
\newcommand{\inc}[1]{\par\smallskip\noindent\textbf{(#1)}\ }
% Resultado destacado
\newcommand{\resultado}[1]{\par\smallskip\noindent\colorbox{caja}{\parbox{\dimexpr\linewidth-2\fboxsep}{#1}}\par\smallskip}
% Idea clave / para entender
\newcommand{\idea}[1]{\par\noindent{\color{tema}\textbf{Idea clave.}} #1\par}
% Código: \codigo{archivo.py}  (se toma de solucionario/codigo/)
\newcommand{\codigo}[1]{\VerbatimInput{\solroot codigo/#1}}

\newcommand{\E}{\mathbf E}
\newcommand{\Hv}{\mathbf H}
\newcommand{\D}{\mathbf D}
\newcommand{\kv}{\mathbf k}
\newcommand{\rv}{\mathbf r}
\newcommand{\J}{\mathbf J}
\newcommand{\T}{\mathbf T}
\newcommand{\R}{\mathbf R}
\newcommand{\Sv}{\mathbf S}
\newcommand{\ud}{\hat{\mathbf u}}
\newcommand{\xh}{\hat{\mathbf x}}
\newcommand{\yh}{\hat{\mathbf y}}
\newcommand{\zh}{\hat{\mathbf z}}
\newcommand{\Ec}{\boldsymbol{\mathcal E}}
\newcommand{\Hc}{\boldsymbol{\mathcal H}}
\newcommand{\Dc}{\boldsymbol{\mathcal D}}
\newcommand{\Bc}{\boldsymbol{\mathcal B}}
\newcommand{\Pc}{\boldsymbol{\mathcal P}}
\newcommand{\cc}{\text{c.c.}}
\newcommand{\dd}{\,\mathrm{d}}
\newcommand{\um}{\si{\micro\metre}}
: "solucionario/" en el maestro, vacío en las pruebas.
\usepackage[spanish,es-noshorthands,es-nodecimaldot,es-tabla]{babel}
\usepackage{fontspec}
\usepackage{amsmath,amssymb,mathtools}
\usepackage{booktabs,array}
\usepackage{xcolor}
\usepackage{enumitem}
\usepackage{graphicx}
\usepackage{float}
\usepackage{siunitx}
\usepackage{fancyvrb}
\usepackage[letterpaper,margin=2.2cm]{geometry}
\usepackage{titlesec}
\usepackage{fancyhdr}
\usepackage{hyperref}
\hypersetup{colorlinks=true,linkcolor=black,urlcolor=black}
\sisetup{output-decimal-marker={.}, per-mode=symbol, range-phrase={--}}
\graphicspath{{\solroot figuras/}}
\fvset{fontsize=\scriptsize, frame=single, numbers=left, numbersep=5pt}

\definecolor{tema}{RGB}{20,60,120}
\definecolor{gris}{RGB}{90,90,90}
\definecolor{caja}{RGB}{235,241,250}
\titleformat{\section}{\Large\bfseries\color{tema}}{\thesection}{0.6em}{}
\titleformat{\subsection}{\large\bfseries\color{tema}}{}{0em}{}
\setlist{itemsep=1pt,topsep=2pt}
\setlength{\parindent}{0pt}
\setlength{\parskip}{0.6ex}
\pagestyle{fancy}\fancyhf{}
\lhead{\small Óptica No Lineal 2026-II — Solucionario del taller de repaso, Examen 1}
\rhead{\small p.~\thepage}

% Encabezado de cada solución: \sol{número}{título}{fuente}
\newcommand{\sol}[3]{\subsection{Problema #1. #2}\noindent{\small\color{gris}#3}\par}
% Inciso: \inc{a}
\newcommand{\inc}[1]{\par\smallskip\noindent\textbf{(#1)}\ }
% Resultado destacado
\newcommand{\resultado}[1]{\par\smallskip\noindent\colorbox{caja}{\parbox{\dimexpr\linewidth-2\fboxsep}{#1}}\par\smallskip}
% Idea clave / para entender
\newcommand{\idea}[1]{\par\noindent{\color{tema}\textbf{Idea clave.}} #1\par}
% Código: \codigo{archivo.py}  (se toma de solucionario/codigo/)
\newcommand{\codigo}[1]{\VerbatimInput{\solroot codigo/#1}}

\newcommand{\E}{\mathbf E}
\newcommand{\Hv}{\mathbf H}
\newcommand{\D}{\mathbf D}
\newcommand{\kv}{\mathbf k}
\newcommand{\rv}{\mathbf r}
\newcommand{\J}{\mathbf J}
\newcommand{\T}{\mathbf T}
\newcommand{\R}{\mathbf R}
\newcommand{\Sv}{\mathbf S}
\newcommand{\ud}{\hat{\mathbf u}}
\newcommand{\xh}{\hat{\mathbf x}}
\newcommand{\yh}{\hat{\mathbf y}}
\newcommand{\zh}{\hat{\mathbf z}}
\newcommand{\Ec}{\boldsymbol{\mathcal E}}
\newcommand{\Hc}{\boldsymbol{\mathcal H}}
\newcommand{\Dc}{\boldsymbol{\mathcal D}}
\newcommand{\Bc}{\boldsymbol{\mathcal B}}
\newcommand{\Pc}{\boldsymbol{\mathcal P}}
\newcommand{\cc}{\text{c.c.}}
\newcommand{\dd}{\,\mathrm{d}}
\newcommand{\um}{\si{\micro\metre}}
: "solucionario/" en el maestro, vacío en las pruebas.
\usepackage[spanish,es-noshorthands,es-nodecimaldot,es-tabla]{babel}
\usepackage{fontspec}
\usepackage{amsmath,amssymb,mathtools}
\usepackage{booktabs,array}
\usepackage{xcolor}
\usepackage{enumitem}
\usepackage{graphicx}
\usepackage{float}
\usepackage{siunitx}
\usepackage{fancyvrb}
\usepackage[letterpaper,margin=2.2cm]{geometry}
\usepackage{titlesec}
\usepackage{fancyhdr}
\usepackage{hyperref}
\hypersetup{colorlinks=true,linkcolor=black,urlcolor=black}
\sisetup{output-decimal-marker={.}, per-mode=symbol, range-phrase={--}}
\graphicspath{{\solroot figuras/}}
\fvset{fontsize=\scriptsize, frame=single, numbers=left, numbersep=5pt}

\definecolor{tema}{RGB}{20,60,120}
\definecolor{gris}{RGB}{90,90,90}
\definecolor{caja}{RGB}{235,241,250}
\titleformat{\section}{\Large\bfseries\color{tema}}{\thesection}{0.6em}{}
\titleformat{\subsection}{\large\bfseries\color{tema}}{}{0em}{}
\setlist{itemsep=1pt,topsep=2pt}
\setlength{\parindent}{0pt}
\setlength{\parskip}{0.6ex}
\pagestyle{fancy}\fancyhf{}
\lhead{\small Óptica No Lineal 2026-II — Solucionario del taller de repaso, Examen 1}
\rhead{\small p.~\thepage}

% Encabezado de cada solución: \sol{número}{título}{fuente}
\newcommand{\sol}[3]{\subsection{Problema #1. #2}\noindent{\small\color{gris}#3}\par}
% Inciso: \inc{a}
\newcommand{\inc}[1]{\par\smallskip\noindent\textbf{(#1)}\ }
% Resultado destacado
\newcommand{\resultado}[1]{\par\smallskip\noindent\colorbox{caja}{\parbox{\dimexpr\linewidth-2\fboxsep}{#1}}\par\smallskip}
% Idea clave / para entender
\newcommand{\idea}[1]{\par\noindent{\color{tema}\textbf{Idea clave.}} #1\par}
% Código: \codigo{archivo.py}  (se toma de solucionario/codigo/)
\newcommand{\codigo}[1]{\VerbatimInput{\solroot codigo/#1}}

\newcommand{\E}{\mathbf E}
\newcommand{\Hv}{\mathbf H}
\newcommand{\D}{\mathbf D}
\newcommand{\kv}{\mathbf k}
\newcommand{\rv}{\mathbf r}
\newcommand{\J}{\mathbf J}
\newcommand{\T}{\mathbf T}
\newcommand{\R}{\mathbf R}
\newcommand{\Sv}{\mathbf S}
\newcommand{\ud}{\hat{\mathbf u}}
\newcommand{\xh}{\hat{\mathbf x}}
\newcommand{\yh}{\hat{\mathbf y}}
\newcommand{\zh}{\hat{\mathbf z}}
\newcommand{\Ec}{\boldsymbol{\mathcal E}}
\newcommand{\Hc}{\boldsymbol{\mathcal H}}
\newcommand{\Dc}{\boldsymbol{\mathcal D}}
\newcommand{\Bc}{\boldsymbol{\mathcal B}}
\newcommand{\Pc}{\boldsymbol{\mathcal P}}
\newcommand{\cc}{\text{c.c.}}
\newcommand{\dd}{\,\mathrm{d}}
\newcommand{\um}{\si{\micro\metre}}
: "solucionario/" en el maestro, vacío en las pruebas.
\usepackage[spanish,es-noshorthands,es-nodecimaldot,es-tabla]{babel}
\usepackage{fontspec}
\usepackage{amsmath,amssymb,mathtools}
\usepackage{booktabs,array}
\usepackage{xcolor}
\usepackage{enumitem}
\usepackage{graphicx}
\usepackage{float}
\usepackage{siunitx}
\usepackage{fancyvrb}
\usepackage[letterpaper,margin=2.2cm]{geometry}
\usepackage{titlesec}
\usepackage{fancyhdr}
\usepackage{hyperref}
\hypersetup{colorlinks=true,linkcolor=black,urlcolor=black}
\sisetup{output-decimal-marker={.}, per-mode=symbol, range-phrase={--}}
\graphicspath{{\solroot figuras/}}
\fvset{fontsize=\scriptsize, frame=single, numbers=left, numbersep=5pt}

\definecolor{tema}{RGB}{20,60,120}
\definecolor{gris}{RGB}{90,90,90}
\definecolor{caja}{RGB}{235,241,250}
\titleformat{\section}{\Large\bfseries\color{tema}}{\thesection}{0.6em}{}
\titleformat{\subsection}{\large\bfseries\color{tema}}{}{0em}{}
\setlist{itemsep=1pt,topsep=2pt}
\setlength{\parindent}{0pt}
\setlength{\parskip}{0.6ex}
\pagestyle{fancy}\fancyhf{}
\lhead{\small Óptica No Lineal 2026-II — Solucionario del taller de repaso, Examen 1}
\rhead{\small p.~\thepage}

% Encabezado de cada solución: \sol{número}{título}{fuente}
\newcommand{\sol}[3]{\subsection{Problema #1. #2}\noindent{\small\color{gris}#3}\par}
% Inciso: \inc{a}
\newcommand{\inc}[1]{\par\smallskip\noindent\textbf{(#1)}\ }
% Resultado destacado
\newcommand{\resultado}[1]{\par\smallskip\noindent\colorbox{caja}{\parbox{\dimexpr\linewidth-2\fboxsep}{#1}}\par\smallskip}
% Idea clave / para entender
\newcommand{\idea}[1]{\par\noindent{\color{tema}\textbf{Idea clave.}} #1\par}
% Código: \codigo{archivo.py}  (se toma de solucionario/codigo/)
\newcommand{\codigo}[1]{\VerbatimInput{\solroot codigo/#1}}

\newcommand{\E}{\mathbf E}
\newcommand{\Hv}{\mathbf H}
\newcommand{\D}{\mathbf D}
\newcommand{\kv}{\mathbf k}
\newcommand{\rv}{\mathbf r}
\newcommand{\J}{\mathbf J}
\newcommand{\T}{\mathbf T}
\newcommand{\R}{\mathbf R}
\newcommand{\Sv}{\mathbf S}
\newcommand{\ud}{\hat{\mathbf u}}
\newcommand{\xh}{\hat{\mathbf x}}
\newcommand{\yh}{\hat{\mathbf y}}
\newcommand{\zh}{\hat{\mathbf z}}
\newcommand{\Ec}{\boldsymbol{\mathcal E}}
\newcommand{\Hc}{\boldsymbol{\mathcal H}}
\newcommand{\Dc}{\boldsymbol{\mathcal D}}
\newcommand{\Bc}{\boldsymbol{\mathcal B}}
\newcommand{\Pc}{\boldsymbol{\mathcal P}}
\newcommand{\cc}{\text{c.c.}}
\newcommand{\dd}{\,\mathrm{d}}
\newcommand{\um}{\si{\micro\metre}}
